% GENERATED FILE. Edit canonical lessons, then run npm run generate:books.
% canonical-source-hash: fnv1a64:8d62ce6ae52c87aa
% canonical-lessons: ML-C154-kaimuttu, ML-C154-kalmuttu, ML-C154-navu, ML-C154-toli, ML-C154-nerri

\chapter{Elbow, Knee, Tongue, Skin, Forehead}
\label{ch:ml-elbow-knee-tongue-skin-forehead}
\hlchaptermodality{154}

\begin{chapteropening}
\textbf{By the end of this chapter:} \emph{I can say an elbow, a knee, a tongue, skin and a forehead.}

Complete the last lesson of chapter 154: i can say an elbow, a knee, a tongue, skin and a forehead.
\end{chapteropening}

\section[\texorpdfstring{\ml{കൈമുട്ട്} (kaimuṭṭŭ)}{കൈമുട്ട് (kaimuṭṭŭ)}]{\ml{കൈമുട്ട്} (kaimuṭṭŭ) — an elbow}

\label{lesson:ML-C154-kaimuttu}



\begin{quote}
Before the new one: say the Malayalam for a mosquito, then the Malayalam for a shoulder.
\end{quote}

\subsection*{You'll want to know: \ml{കൈമുട്ട്}}
\textbf{\ml{കൈമുട്ട്}} — \emph{kaimuṭṭŭ} — \textquotedblleft{}an elbow\textquotedblright{}.

\begin{grammarlens}[title={What you've built}]
One more word for this chapter.
\end{grammarlens}

\subsection*{Your turn}
\begin{itemize}
\raggedright
  \item \emph{Say it:} \emph{kaimuṭṭŭ}
  \item \emph{Say it:} \emph{kaimuṭṭŭ}, once more
  \item \emph{Say it:} say \emph{kaimuṭṭŭ}
  \item \emph{Recall:} say the Malayalam for a mosquito, then the Malayalam for a shoulder, then say \emph{kaimuṭṭŭ} again
\end{itemize}

\begin{tcolorbox}[breakable,colback=teal!4,colframe=teal!35!black,title={Before you move on}]
What does \ml{കൈമുട്ട്} mean? (\textquotedblleft{}An elbow\textquotedblright{}.) Say it once more.
\end{tcolorbox}
\section[\texorpdfstring{\ml{കാൽമുട്ട്} (kālmuṭṭŭ)}{കാൽമുട്ട് (kālmuṭṭŭ)}]{\ml{കാൽമുട്ട്} (kālmuṭṭŭ) — a knee}

\label{lesson:ML-C154-kalmuttu}



\begin{quote}
Before the new one: say the Malayalam for a shoulder, then the Malayalam for an elbow.
\end{quote}

\subsection*{You'll want to know: \ml{കാൽമുട്ട്}}
\textbf{\ml{കാൽമുട്ട്}} — \emph{kālmuṭṭŭ} — \textquotedblleft{}a knee\textquotedblright{}.

\begin{grammarlens}[title={What you've built}]
One more word for this chapter.
\end{grammarlens}

\subsection*{Your turn}
\begin{itemize}
\raggedright
  \item \emph{Say it:} \emph{kālmuṭṭŭ}
  \item \emph{Say it:} \emph{kālmuṭṭŭ}, once more
  \item \emph{Say it:} say \emph{kālmuṭṭŭ}
  \item \emph{Recall:} say the Malayalam for a shoulder, then the Malayalam for an elbow, then say \emph{kālmuṭṭŭ} again
\end{itemize}

\begin{tcolorbox}[breakable,colback=teal!4,colframe=teal!35!black,title={Before you move on}]
What does \ml{കാൽമുട്ട്} mean? (\textquotedblleft{}A knee\textquotedblright{}.) Say it once more.
\end{tcolorbox}
\section[\texorpdfstring{\ml{നാവ്} (nāvŭ)}{നാവ് (nāvŭ)}]{\ml{നാവ്} (nāvŭ) — a tongue}

\label{lesson:ML-C154-navu}



\begin{quote}
Before the new one: say the Malayalam for an elbow, then the Malayalam for a knee.

Say \textquotedblleft{}I like Malayalam.\textquotedblright{} (\emph{Enikkŭ malayāḷaṁ iṣṭamāṇŭ} — \textquotedblleft{}to me, Malayalam is a liking.\textquotedblright{})
\end{quote}

\subsection*{You'll want to know: \ml{നാവ്}}
\textbf{\ml{നാവ്}} — \emph{nāvŭ} — \textquotedblleft{}a tongue\textquotedblright{}.

\begin{grammarlens}[title={What you've built}]
One more word for this chapter.
\end{grammarlens}

\subsection*{Your turn}
\begin{itemize}
\raggedright
  \item \emph{Say it:} \emph{nāvŭ}
  \item \emph{Say it:} \emph{nāvŭ}, once more
  \item \emph{Say it:} say \emph{nāvŭ}
  \item \emph{Recall:} say the Malayalam for an elbow, then the Malayalam for a knee, then say \emph{nāvŭ} again
\end{itemize}

\begin{tcolorbox}[breakable,colback=teal!4,colframe=teal!35!black,title={Before you move on}]
What does \ml{നാവ്} mean? (\textquotedblleft{}A tongue\textquotedblright{}.) Say it once more.
\end{tcolorbox}
\section[\texorpdfstring{\ml{തൊലി} (tolī)}{തൊലി (tolī)}]{\ml{തൊലി} (tolī) — skin}

\label{lesson:ML-C154-toli}



\begin{quote}
Before the new one: say the Malayalam for a knee, then the Malayalam for a tongue.
\end{quote}

\subsection*{You'll want to know: \ml{തൊലി}}
\textbf{\ml{തൊലി}} — \emph{tolī} — \textquotedblleft{}skin\textquotedblright{}.

\begin{grammarlens}[title={What you've built}]
One more word for this chapter.
\end{grammarlens}

\subsection*{Your turn}
\begin{itemize}
\raggedright
  \item \emph{Say it:} \emph{tolī}
  \item \emph{Say it:} \emph{tolī}, once more
  \item \emph{Say it:} say \emph{tolī}
  \item \emph{Recall:} say the Malayalam for a knee, then the Malayalam for a tongue, then say \emph{tolī} again
\end{itemize}

\begin{tcolorbox}[breakable,colback=teal!4,colframe=teal!35!black,title={Before you move on}]
What does \ml{തൊലി} mean? (\textquotedblleft{}Skin\textquotedblright{}.) Say it once more.
\end{tcolorbox}
\section[\texorpdfstring{\ml{നെറ്റി} (neṟṟi)}{നെറ്റി (neṟṟi)}]{\ml{നെറ്റി} (neṟṟi) — a forehead}

\label{lesson:ML-C154-nerri}



\begin{quote}
Before the new one: say the Malayalam for a tongue, then the Malayalam for skin.

In \emph{enikkŭ malayāḷaṁ iṣṭamāṇŭ}, which word names the one who likes? (\emph{Enikkŭ}, \textquotedblleft{}to me.\textquotedblright{})
\end{quote}

\subsection*{You'll want to know: \ml{നെറ്റി}}
\textbf{\ml{നെറ്റി}} — \emph{neṟṟi} — \textquotedblleft{}a forehead\textquotedblright{}.

\begin{grammarlens}[title={What you've built}]
One more word for this chapter.
\end{grammarlens}

\subsection*{Your turn}
\begin{itemize}
\raggedright
  \item \emph{Say it:} \emph{neṟṟi}
  \item \emph{Say it:} \emph{neṟṟi}, once more
  \item \emph{Say it:} say \emph{neṟṟi}
  \item \emph{Recall:} say the Malayalam for a tongue, then the Malayalam for skin, then say \emph{neṟṟi} again
\end{itemize}

\begin{tcolorbox}[breakable,colback=teal!4,colframe=teal!35!black,title={Before you move on}]
What does \ml{നെറ്റി} mean? (\textquotedblleft{}A forehead\textquotedblright{}.) Say it once more.
\end{tcolorbox}
